% !PRINTED 265
\weberpage{265}
%<<CONT
worin die $M_1, M_2 \ldots M_k$ rational von den $a_{i,k}, b_{i,k}$ abhängen. Insbesondere ist
\[
M_k = \Sigma \pm b_{1,1} b_{2,2} \ldots b_{k,k},
\]
und man kann die $b_{i,k}$ immer so annehmen, dass $M_k$ von Null verschieden ist. Nach \S.\,31 kann man also $\varepsilon$ so annehmen, dass, wenn $R_k$ von Null verschieden ist, $R_k'$ dasselbe Zeichen hat, wie $R_k$, und wenn $R_k$ verschwindet, $R_k'$ nicht verschwindet.

Wir können das hierdurch Bewiesene mit den Ergebnissen des \S.\,81 in eine allgemeine Regel zur Bestimmung der Anzahl der negativen Quadrate, auch für den Fall verschwindender Determinante, zusammenfassen.

% The source continues to use the same variant rho (varrho) introduced in section 80.
VI. \gesperrt{Um die Zahlen $\pi$, $\nu$, $\varrho$ der Function $\varphi(x)$ zu bestimmen, ordne man die Variablen $x_1, x_2 \ldots x_m$ so an, dass in der Kette der Haupt-Unterdeterminanten}
\weq{7}{R_m, R_{m-1} \ldots R_1, R_0}
\gesperrt{eine möglichst kleine Anzahl von Anfangsgliedern verschwindet, und dass von den folgenden Gliedern nicht zwei neben einander stehende verschwinden; $\varrho$ ist dann die Anzahl der verschwindenden Anfangsglieder, $\nu$ die Anzahl der Zeichenwechsel und $\pi = m - \nu - \varrho$.}

Schliesslich sei noch bemerkt, dass man die Anzahl der Zeichenwechsel in der Reihe (7) auch von rechts nach links abzählen kann, d. h. dass man dieselbe Anzahl von Zeichenwechseln findet, wenn man die Reihe in umgekehrter Ordnung schreibt.

\begin{center}
\S.\,84.

\gesperrt{Anwendung auf die Bezoutiante.}
\end{center}

Die genaue Discussion der Trägheit der quadratischen Formen hatte für uns nur den Zweck, die Anzahl der reellen und imaginären Wurzeln einer algebraischen Gleichung durch Abzählung der positiven und negativen Quadrate der Bezoutiante zu bestimmen. Ist $n$ der Grad der Gleichung, so ist ihre Bezoutiante eine quadratische Form von $n - 1$ Veränderlichen, und im \S.\,73 haben wir sie für $n = 3, 4, 5$ vollständig gebildet
%CONT>>
